% Generated draft: requires mathematical and source review.
\documentclass[11pt]{article}
\usepackage[margin=25mm]{geometry}
\usepackage{fontspec}
\setmainfont{Latin Modern Roman}
\usepackage{amsmath,amssymb,mathtools,unicode-math}
\setmathfont{Latin Modern Math}
\usepackage{enumitem,needspace,xurl,hyperref}
\setlength{\parindent}{0pt}
\setlength{\parskip}{5pt}
\setlength{\emergencystretch}{3em}
\newcommand{\N}{\mathbb{N}}
\newcommand{\Z}{\mathbb{Z}}
\newcommand{\Q}{\mathbb{Q}}
\newcommand{\R}{\mathbb{R}}
\newcommand{\C}{\mathbb{C}}
\newcommand{\F}{\mathbb{F}}
\newcommand{\A}{\mathbb{A}}
\newcommand{\PP}{\mathbb{P}}
\newcommand{\E}{\mathbb{E}}
\newcommand{\eps}{\varepsilon}
\newcommand{\dd}{\,\mathrm{d}}
\begin{document}
\section*{UnsolvedMath ID: 30006063}
\textbf{Source problem code:} catalogue entry
\subsection{Definitions and mathematical statement}
The supplied record asks about coprime integers \(d,r\) with \(r>0\) and a fixed Feigin--Odesskii Poisson bracket of type \((d,r)\). The notation \(\mathbb P^{d-1}\) requires \(d\geq 1\). The record does not specify the base field, the complete admissibility conditions for this type, or the construction and parameters of the fixed bracket.

The following makes the question precise under the explicitly provisional interpretation that the brackets are regular algebraic Poisson brackets. Let \(k\) be the intended base field, let \((d,r)\) be an admissible pair satisfying the supplied conditions, and put \(X=\mathbb P^{d-1}_k\). Suppose that the specified Feigin--Odesskii bracket \(\pi_0\) of type \((d,r)\) on \(X\) has been fixed. Write
\[
 B(X)=H^0\!\left(X,\mathcal Hom_{\mathcal O_X}\!\left(\bigwedge^2\Omega^1_{X/k},\mathcal O_X\right)\right).
\]
An element \(\pi\in B(X)\) gives an alternating \(k\)-bilinear bracket \(\{f,g\}_{\pi}\) on local regular functions, satisfying the Leibniz rule in each argument. It is a Poisson bracket if, on every open set and for all local regular functions \(f,g,h\),
\[
 \{f,\{g,h\}_{\pi}\}_{\pi}
 +\{g,\{h,f\}_{\pi}\}_{\pi}
 +\{h,\{f,g\}_{\pi}\}_{\pi}=0.
\]
A \(k\)-linear subspace \(V\subseteq B(X)\) is a space of compatible Poisson brackets if every element of \(V\) is Poisson; equivalently, every finite \(k\)-linear combination of elements of \(V\) is Poisson.

Determine
\[
 M_k(X,\pi_0)=\max\left\{\dim_k V:
 V\subseteq B(X)\text{ is a }k\text{-linear subspace},\quad
 \pi_0\in V,\quad
 \text{every }\pi\in V\text{ is Poisson}\right\}.
\]
The question is asked for each fixed bracket \(\pi_0\); no independence from its parameters is assumed. The base field, admissible type conventions, and interpretation as regular algebraic brackets must be confirmed before this conditional formulation can be treated as a definitive statement of the source problem.
\subsection{Short English statement}
Given a fixed Feigin--Odesskii bracket of type \((d,r)\), with \(r>0\) and \(\gcd(d,r)=1\), determine the largest dimension of a vector space of Poisson brackets on \(\mathbb P^{d-1}\) containing it, such that every linear combination is again Poisson. The base field, admissible types, and class of brackets remain source qualifications requiring review.
\subsection{Sources}
\begin{itemize}
\item Supplied local UnsolvedMath registry record: permanent ID \(30006063\), title \emph{Compatible Feigin--Odesskii Poisson Brackets}, source problem code \texttt{catalogue entry}, in the Oberwolfach Reports Open Problems collection. No individual proposer is supplied; the proposed year is recorded as 2024.
\item Supplied report citation: \emph{MFO-RIMS Tandem Workshop: Algebraic Geometry and Noncommutative Projective Varieties (2025)}, Oberwolfach Reports, pp.\ 2551--2606, DOI \texttt{10.4171/owr/2024/44}. \url{https://doi.org/10.4171/owr/2024/44}. The original excerpt carries page number 2571. The supplied literature list identifies this source as Oberwolfach Report 44/2024; the citation's 2025 designation and the record's 2024 dating are retained without reconciliation.
\item Supplied literature reference: Markarian--Polishchuk, \emph{Compatible Feigin--Odesskii Poisson brackets}. \url{https://doi.org/10.1007/s00229-023-01479-2}.
\item Literature assessment dated exactly \texttt{2026-08-22}: the supplied curation labels the problem open and states that the general maximal-dimension classification remains unresolved, while special cases determine dimensions such as 9 or 10. The hypotheses of those special cases are not supplied. This is supplied dated curation derived from the Oberwolfach material, not an independently verified literature assessment; its status and cited claims require independent checking.
\item The registry labels the statement \texttt{corrected\_verified} and reports that it was reconstructed and checked against the official Oberwolfach report. This is a supplied verification claim, not verification performed for this draft.
\end{itemize}
\end{document}
